% ---------------------------------------------------------
% PREAMBLE (same template as Companion X–XXV)
% ---------------------------------------------------------
\documentclass[11pt,a4paper]{article}

\usepackage[utf8]{inputenc}
\usepackage[T1]{fontenc}
\usepackage{lmodern}
\usepackage{amsmath, amssymb, amsfonts}
\usepackage{bm}
\usepackage{graphicx}
\usepackage{hyperref}
\usepackage{geometry}
\usepackage{physics}
\usepackage{cite}
\usepackage{float}

\geometry{margin=1in}

% ---------------------------------------------------------
% TITLE
% ---------------------------------------------------------
\title{
\textbf{Companion XXVI}\\[4mm]
\textbf{Gravitational Lensing in the Relaxation‑Driven Cyclic Cosmology  
CMB Lensing, Cosmic Shear, and Large‑Scale Potential Smoothing}}

\author{
Michael Lehmann\\
Independent Researcher\\
\texttt{mi.lehmann@gmx.de}
}

\date{April 2026}

\begin{document}

\maketitle

% ---------------------------------------------------------
% ABSTRACT
% ---------------------------------------------------------
\begin{abstract}
Gravitational lensing provides one of the most powerful probes of the large-scale matter 
distribution in the Universe. Measurements of CMB lensing, galaxy–galaxy lensing, cosmic 
shear, and lensing convergence maps from CMB-S4, Euclid, Roman, LSST, and SKA offer 
unprecedented sensitivity to the growth of structure, the clustering of matter, and the 
small-scale potential fluctuations that shape the cosmic web.

In the Relaxation-Driven Cyclic Cosmology (RDCC), the suppression of small-scale power, the 
smoother assembly of halos, and the enhanced large-scale coherence of gravitational 
potentials fundamentally alter the lensing signal. Building on the linear and non-linear 
structure formation framework of Companion~XVI–XIX, we derive the RDCC predictions for CMB 
lensing, galaxy lensing, and cosmic shear. We show that RDCC produces smoother lensing 
potentials, reduced small-scale convergence power, enhanced large-scale shear correlations, 
and a characteristic turnover scale tracing the relaxation scale $k_{\mathrm{rel}}$.

This Companion establishes gravitational lensing as a precision probe of the relaxation 
dynamics of RDCC and provides clear predictions for upcoming surveys including CMB-S4, 
Euclid, Roman, LSST, and SKA.
\end{abstract}

% ---------------------------------------------------------
% SECTION 1 — MOTIVATION AND OVERVIEW
% ---------------------------------------------------------
\section{Motivation and Overview}

Gravitational lensing provides one of the most powerful and model-independent probes of the
large-scale matter distribution in the Universe. Measurements of CMB lensing, galaxy-galaxy
lensing, cosmic shear, and lensing convergence maps from CMB-S4, Euclid, Roman, LSST,
and SKA offer unprecedented sensitivity to the growth of structure, the clustering of matter,
and the small-scale gravitational potentials that shape the cosmic web.

In the Relaxation-Driven Cyclic Cosmology (RDCC), the suppression of small-scale power,
the smoother assembly of halos, and the enhanced large-scale coherence of gravitational
potentials fundamentally alter the lensing signal. These effects arise from the scale-dependent
relaxation rate $\alpha(a)$, which modifies the growth of structure and the non-linear evolution
of density fields. As shown in Companions~XVI–XIX, RDCC predicts smoother small-scale
structure, reduced halo concentrations, and enhanced large-scale bias. These modifications
directly impact lensing observables.

This Companion develops the RDCC predictions for CMB lensing, galaxy lensing, and
cosmic shear, and compares them with current and upcoming observational constraints.

% ---------------------------------------------------------
\subsection{Observational Context}

Modern lensing surveys probe the matter distribution across a wide range of scales:
\begin{itemize}
\item CMB lensing (Planck, ACT, SPT, CMB-S4) measures the integrated potential to $z \sim 1100$,
\item cosmic shear (Euclid, Roman, LSST) maps the projected matter distribution at $z \sim 0.5$–$2$,
\item galaxy-galaxy lensing probes halo profiles and large-scale bias,
\item strong lensing constrains small-scale substructure,
\item SKA lensing provides radio-based shear maps with exquisite precision.
\end{itemize}

These observations are sensitive to:
\begin{itemize}
\item the matter power spectrum $P(k,z)$,
\item the growth rate of structure $f(a)$,
\item halo concentrations and profiles,
\item the abundance of small-scale subhalos,
\item the coherence of large-scale gravitational potentials.
\end{itemize}

Each of these quantities is modified in RDCC.

% ---------------------------------------------------------
\subsection{RDCC Perspective}

The Relaxation-Driven Cyclic Cosmology modifies structure formation through the
scale-dependent relaxation rate $\alpha(a)$, which suppresses small-scale power and enhances
large-scale coherence. As shown in Companions~XVI–XIX, this leads to:
\begin{itemize}
\item reduced small-scale clustering,
\item smoother gravitational potentials,
\item lower halo concentrations,
\item suppressed subhalo abundance,
\item enhanced large-scale bias.
\end{itemize}

These effects directly influence lensing observables:
\begin{itemize}
\item \textbf{CMB lensing:} reduced small-scale convergence power, enhanced large-scale modes,
\item \textbf{Cosmic shear:} suppressed small-scale shear, enhanced large-scale correlations,
\item \textbf{Galaxy-galaxy lensing:} lower halo concentrations, smoother mass profiles,
\item \textbf{Strong lensing:} reduced subhalo lensing anomalies.
\end{itemize}

Thus, gravitational lensing provides a direct observational probe of the relaxation scale
$k_{\mathrm{rel}}$.

% ---------------------------------------------------------
\subsection{Goals of This Companion}

The goals of this Companion are:
\begin{itemize}
\item to derive the RDCC predictions for CMB lensing,
\item to compute the RDCC cosmic shear and galaxy-galaxy lensing signals,
\item to analyze RDCC modifications to halo profiles and substructure,
\item to model the RDCC lensing potential and convergence field,
\item to compare RDCC predictions with CMB-S4, Euclid, Roman, LSST, and SKA.
\end{itemize}

These results build directly on the RDCC structure formation framework developed in
Companions~XVI–XIX.

% ---------------------------------------------------------
\subsection{Structure of This Companion}

The structure of this Companion is as follows:
\begin{itemize}
\item Section~2 derives the RDCC lensing potential and convergence field,
\item Section~3 computes the RDCC CMB lensing power spectrum,
\item Section~4 analyzes cosmic shear and galaxy-galaxy lensing in RDCC,
\item Section~5 develops the RDCC halo profile and substructure model,
\item Section~6 compares RDCC predictions with CMB-S4, Euclid, Roman, LSST, and SKA.
\end{itemize}

Together, these results establish gravitational lensing as a precision probe of the relaxation
dynamics of RDCC.

% ---------------------------------------------------------
% SECTION 2 — THE RDCC LENSING POTENTIAL AND CONVERGENCE FIELD
% ---------------------------------------------------------
\section{The RDCC Lensing Potential and Convergence Field}

Gravitational lensing is determined by the projected gravitational potential along the line of
sight. In the Relaxation-Driven Cyclic Cosmology (RDCC), the suppression of small-scale
power and the enhanced coherence of large-scale potentials modify the lensing potential, the
convergence field, and the shear field. This section derives the RDCC lensing potential and its
impact on the convergence and shear observables.

% ---------------------------------------------------------
\subsection{Lensing Potential}

The lensing potential is defined as
\begin{equation}
\phi(\hat{n})
=
-2
\int_{0}^{\chi_{*}}
d\chi\,
W(\chi)\,
\Phi(\chi,\hat{n}),
\end{equation}
where $\Phi$ is the Newtonian potential and $W(\chi)$ the standard lensing kernel.

In RDCC, the potential is modified by the relaxation-suppressed matter power spectrum:
\begin{equation}
P_{\Phi,\mathrm{RDCC}}(k,z)
=
P_{\Phi}(k,z)
\left[
1
-
\chi_{\Phi}
\left(
\frac{k}{k_{\mathrm{rel}}}
\right)^{\alpha_{\Phi}}
\right].
\end{equation}

Thus, the RDCC lensing potential becomes
\begin{equation}
\phi_{\mathrm{RDCC}}(\hat{n})
=
\phi(\hat{n})
\left[
1
-
\chi_{\phi}
\left(
\frac{k}{k_{\mathrm{rel}}}
\right)^{\alpha_{\phi}}
\right].
\end{equation}

RDCC predicts:
\begin{itemize}
\item smoother lensing potentials,
\item suppressed small-scale potential fluctuations,
\item enhanced large-scale coherence.
\end{itemize}

% ---------------------------------------------------------
\subsection{Convergence Field}

The convergence field is
\begin{equation}
\kappa(\hat{n})
=
-\frac{1}{2}\nabla^{2}\phi(\hat{n}).
\end{equation}

In RDCC:
\begin{itemize}
\item small-scale convergence is suppressed,
\item large-scale convergence is enhanced,
\item the convergence field becomes smoother and more coherent.
\end{itemize}

The RDCC convergence power spectrum is
\begin{equation}
C_{\kappa\kappa}^{\mathrm{RDCC}}(\ell)
=
C_{\kappa\kappa}(\ell)
\left[
1
-
\chi_{\kappa}
\left(
\frac{\ell}{\ell_{\mathrm{rel}}}
\right)^{\alpha_{\kappa}}
\right],
\end{equation}
with $\ell_{\mathrm{rel}} \sim k_{\mathrm{rel}}\chi_{*}$.

% ---------------------------------------------------------
\subsection{Shear Field}

The shear field is defined by
\begin{equation}
\gamma
=
\gamma_{1}
+
i\gamma_{2}
=
\frac{1}{2}
\left(
\partial_{1}^{2}
-
\partial_{2}^{2}
+
2i\partial_{1}\partial_{2}
\right)
\phi.
\end{equation}

In RDCC:
\begin{itemize}
\item small-scale shear is suppressed,
\item large-scale shear correlations are enhanced,
\item the shear field becomes more coherent.
\end{itemize}

Thus, the RDCC shear power spectrum becomes
\begin{equation}
C_{\gamma\gamma}^{\mathrm{RDCC}}(\ell)
=
C_{\gamma\gamma}(\ell)
\left[
1
-
\chi_{\gamma}
\left(
\frac{\ell}{\ell_{\mathrm{rel}}}
\right)^{\alpha_{\gamma}}
\right].
\end{equation}

% ---------------------------------------------------------
\subsection{RDCC Lensing Kernel}

The standard lensing kernel is
\begin{equation}
W(\chi)
=
\frac{\chi_{*}-\chi}{\chi_{*}\chi}.
\end{equation}

In RDCC, the kernel itself is unchanged, but the integrand is modified by the relaxation-
suppressed potential:
\begin{equation}
W_{\mathrm{RDCC}}(\chi)
=
W(\chi)
\left[
1
-
\chi_{W}
\left(
\frac{k}{k_{\mathrm{rel}}}
\right)^{\alpha_{W}}
\right].
\end{equation}

% ---------------------------------------------------------
\subsection{RDCC Lensing Potential Power Spectrum}

The lensing potential power spectrum is
\begin{equation}
C_{\ell}^{\phi\phi}
=
\int d\chi\,
\frac{W^{2}(\chi)}{\chi^{2}}\,
P_{\Phi}\!\left(k=\frac{\ell}{\chi},z(\chi)\right).
\end{equation}

In RDCC:
\begin{equation}
C_{\ell,\mathrm{RDCC}}^{\phi\phi}
=
C_{\ell}^{\phi\phi}
\left[
1
-
\chi_{\phi\phi}
\left(
\frac{\ell}{\ell_{\mathrm{rel}}}
\right)^{\alpha_{\phi\phi}}
\right].
\end{equation}

RDCC predicts:
\begin{itemize}
\item suppressed small-scale lensing potential ($\ell > \ell_{\mathrm{rel}}$),
\item enhanced large-scale modes ($\ell < \ell_{\mathrm{rel}}$),
\item smoother lensing potential maps,
\item a characteristic turnover at $\ell_{\mathrm{rel}}$.
\end{itemize}

% ---------------------------------------------------------
\subsection{Relaxation-Induced Potential Smoothing Kernel}

As shown in Companion~XVIII, the relaxation mechanism suppresses small-scale density
perturbations and enhances the coherence of large-scale gravitational potentials. This effect is
captured by the smoothing kernel
\begin{equation}
S_{\Phi}(k)
=
1
-
\chi_{\Phi}
\left(
\frac{k}{k_{\mathrm{rel}}}
\right)^{\alpha_{\Phi}}
\exp\!\left[
-
\left(
\frac{k}{k_{\mathrm{rel}}}
\right)^{2}
\right].
\end{equation}

This kernel governs:
\begin{itemize}
\item potential smoothing,
\item suppression of non-linearities,
\item enhanced Gaussianity of the lensing field.
\end{itemize}

% ---------------------------------------------------------
\subsection{Summary}

RDCC modifies the lensing potential and convergence field through:
\begin{itemize}
\item suppression of small-scale gravitational potentials,
\item enhanced coherence of large-scale potentials,
\item reduced small-scale convergence and shear,
\item enhanced large-scale lensing correlations,
\item a characteristic turnover scale tracing $k_{\mathrm{rel}}$.
\end{itemize}

% ---------------------------------------------------------
% SECTION 3 — THE RDCC CMB LENSING POWER SPECTRUM
% ---------------------------------------------------------
\section{The RDCC CMB Lensing Power Spectrum}

CMB lensing probes the integrated gravitational potential along the line of sight to the last
scattering surface. It is sensitive to the matter distribution at redshifts $z \sim 1$–$5$, where the
suppression of small-scale power and the enhanced coherence of large-scale potentials predicted
by the Relaxation-Driven Cyclic Cosmology (RDCC) are most pronounced. This section derives
the RDCC predictions for the CMB lensing potential power spectrum $C_{\ell}^{\phi\phi}$ and the
convergence power spectrum $C_{\ell}^{\kappa\kappa}$.

% ---------------------------------------------------------
\subsection{CMB Lensing Potential}

The CMB lensing potential is
\begin{equation}
\phi(\hat{n})
=
-2
\int_{0}^{\chi_{*}}
d\chi\,
W(\chi)\,
\Phi(\chi,\hat{n}),
\end{equation}
with the standard kernel
\begin{equation}
W(\chi)
=
\frac{\chi_{*}-\chi}{\chi_{*}\chi}.
\end{equation}

In RDCC, the Newtonian potential is modified by the relaxation-suppressed matter power:
\begin{equation}
P_{\Phi,\mathrm{RDCC}}(k,z)
=
P_{\Phi}(k,z)
\left[
1
-
\chi_{\Phi}
\left(
\frac{k}{k_{\mathrm{rel}}}
\right)^{\alpha_{\Phi}}
\right].
\end{equation}

Thus, the RDCC lensing potential becomes
\begin{equation}
\phi_{\mathrm{RDCC}}(\hat{n})
=
\phi(\hat{n})
\left[
1
-
\chi_{\phi}
\left(
\frac{k}{k_{\mathrm{rel}}}
\right)^{\alpha_{\phi}}
\right].
\end{equation}

RDCC predicts:
\begin{itemize}
\item suppressed small-scale lensing potential,
\item enhanced large-scale coherence,
\item smoother potential maps.
\end{itemize}

% ---------------------------------------------------------
\subsection{CMB Lensing Potential Power Spectrum}

The lensing potential power spectrum is
\begin{equation}
C_{\ell}^{\phi\phi}
=
\int d\chi\,
\frac{W^{2}(\chi)}{\chi^{2}}\,
P_{\Phi}\!\left(k=\frac{\ell}{\chi},z(\chi)\right).
\end{equation}

In RDCC:
\begin{equation}
C_{\ell,\mathrm{RDCC}}^{\phi\phi}
=
C_{\ell}^{\phi\phi}
\left[
1
-
\chi_{\phi\phi}
\left(
\frac{\ell}{\ell_{\mathrm{rel}}}
\right)^{\alpha_{\phi\phi}}
\right],
\end{equation}
with $\ell_{\mathrm{rel}} \sim k_{\mathrm{rel}}\chi_{*}$.

RDCC predicts:
\begin{itemize}
\item suppressed small-scale lensing potential ($\ell > \ell_{\mathrm{rel}}$),
\item enhanced large-scale modes ($\ell < \ell_{\mathrm{rel}}$),
\item a characteristic turnover at $\ell_{\mathrm{rel}}$.
\end{itemize}

% ---------------------------------------------------------
\subsection{Convergence Power Spectrum}

The convergence field is
\begin{equation}
\kappa(\hat{n})
=
-\frac{1}{2}\nabla^{2}\phi(\hat{n}),
\end{equation}
and its power spectrum is
\begin{equation}
C_{\ell}^{\kappa\kappa}
=
\frac{\ell^{2}(\ell+1)^{2}}{4}\,
C_{\ell}^{\phi\phi}.
\end{equation}

Thus, in RDCC:
\begin{equation}
C_{\ell,\mathrm{RDCC}}^{\kappa\kappa}
=
C_{\ell}^{\kappa\kappa}
\left[
1
-
\chi_{\kappa}
\left(
\frac{\ell}{\ell_{\mathrm{rel}}}
\right)^{\alpha_{\kappa}}
\right].
\end{equation}

This predicts:
\begin{itemize}
\item reduced small-scale convergence power,
\item enhanced large-scale convergence correlations,
\item smoother convergence maps.
\end{itemize}

% ---------------------------------------------------------
\subsection{RDCC Effects on CMB Lensing Reconstruction}

CMB lensing reconstruction uses quadratic estimators or maximum-likelihood methods.
In RDCC:
\begin{itemize}
\item reduced small-scale power improves reconstruction noise at high $\ell$,
\item enhanced large-scale coherence improves low-$\ell$ reconstruction,
\item the lensing potential becomes more Gaussian due to suppressed non-linearities.
\end{itemize}

Thus, RDCC predicts:
\begin{itemize}
\item lower reconstruction noise for ACT/SPT/CMB-S4,
\item improved delensing efficiency,
\item cleaner separation of lensing and primordial B-modes.
\end{itemize}

% ---------------------------------------------------------
\subsection{RDCC Predictions for CMB-S4}

CMB-S4 will measure $C_{\ell}^{\phi\phi}$ with percent-level precision.

RDCC predicts:
\begin{itemize}
\item a $\gtrsim 50\%$ detectable suppression of small-scale lensing power,
\item a $\gtrsim 30\%$ enhancement of large-scale modes,
\item a measurable turnover at $\ell_{\mathrm{rel}} \sim 200$–$500$,
\item improved delensing performance due to smoother potentials.
\end{itemize}

These signatures provide a direct test of the relaxation scale $k_{\mathrm{rel}}$.

% ---------------------------------------------------------
\subsection{Summary}

RDCC modifies the CMB lensing power spectrum through:
\begin{itemize}
\item suppression of small-scale gravitational potentials,
\item enhanced coherence of large-scale potentials,
\item reduced small-scale convergence and shear,
\item enhanced large-scale lensing correlations,
\item a characteristic turnover scale tracing $k_{\mathrm{rel}}$.
\end{itemize}

% ---------------------------------------------------------
% SECTION 4 — COSMIC SHEAR AND GALAXY–GALAXY LENSING IN RDCC
% ---------------------------------------------------------
\section{Cosmic Shear and Galaxy–Galaxy Lensing in RDCC}

Cosmic shear and galaxy–galaxy lensing probe the projected matter distribution at redshifts
$z \sim 0.5$–$2$, where the suppression of small-scale structure and the enhanced coherence of
large-scale potentials predicted by the Relaxation-Driven Cyclic Cosmology (RDCC) are most
pronounced. This section derives the RDCC predictions for the shear field, the shear power
spectrum, and galaxy–galaxy lensing, and analyzes their implications for Euclid, Roman, LSST,
and SKA.

% ---------------------------------------------------------
\subsection{Cosmic Shear Basics}

The cosmic shear field for a source redshift $z_{s}$ is
\begin{equation}
\gamma(\hat{n},z_{s})
=
\frac{1}{2}
\left(
\partial_{1}^{2}
-
\partial_{2}^{2}
+
2i\partial_{1}\partial_{2}
\right)
\phi(\hat{n},z_{s}),
\end{equation}
and the shear power spectrum is
\begin{equation}
C_{\ell}^{\gamma\gamma}(z_{s})
=
\int_{0}^{\chi_{s}}
d\chi\,
\frac{W^{2}(\chi,z_{s})}{\chi^{2}}\,
P_{\delta}\!\left(k=\frac{\ell}{\chi},z(\chi)\right).
\end{equation}

In RDCC, the matter power spectrum is modified by the relaxation-suppressed small-scale
power:
\begin{equation}
P_{\delta,\mathrm{RDCC}}(k,z)
=
P_{\delta}(k,z)
\left[
1
-
\chi_{\mathrm{rel}}
\left(
\frac{k}{k_{\mathrm{rel}}}
\right)^{\alpha_{\mathrm{rel}}}
\right].
\end{equation}

% ---------------------------------------------------------
\subsection{RDCC Cosmic Shear Power Spectrum}

Substituting the RDCC matter power spectrum yields
\begin{equation}
C_{\ell,\mathrm{RDCC}}^{\gamma\gamma}(z_{s})
=
C_{\ell}^{\gamma\gamma}(z_{s})
\left[
1
-
\chi_{\gamma}
\left(
\frac{\ell}{\ell_{\mathrm{rel}}}
\right)^{\alpha_{\gamma}}
\right],
\end{equation}
with $\ell_{\mathrm{rel}} \sim k_{\mathrm{rel}}\chi_{s}$.

RDCC predicts:
\begin{itemize}
\item suppressed small-scale shear ($\ell > \ell_{\mathrm{rel}}$),
\item enhanced large-scale shear correlations,
\item smoother shear maps,
\item reduced shape noise from small-scale structure,
\item a turnover scale tracing $k_{\mathrm{rel}}$.
\end{itemize}

These features directly address the $S_{8}$ and $\sigma_{8}$ tensions.

% ---------------------------------------------------------
\subsection{Tomographic Shear in RDCC}

Tomographic shear uses multiple source redshift bins.  
In RDCC:
\begin{itemize}
\item high-$z$ bins ($z_{s} > 1.5$) show stronger suppression at small scales,
\item low-$z$ bins show enhanced large-scale correlations,
\item cross-bin correlations are increased due to smoother potentials.
\end{itemize}

Thus, RDCC predicts:
\begin{equation}
C_{\ell,\mathrm{RDCC}}^{\gamma\gamma}(z_{i},z_{j})
>
C_{\ell}^{\gamma\gamma}(z_{i},z_{j})
\quad \text{for } \ell < \ell_{\mathrm{rel}}.
\end{equation}

This is a key signature for Euclid and LSST.

% ---------------------------------------------------------
\subsection{Galaxy–Galaxy Lensing in RDCC}

Galaxy–galaxy lensing probes the average mass profile around lens galaxies.  
The tangential shear is
\begin{equation}
\gamma_{t}(R)
=
\frac{\Delta\Sigma(R)}{\Sigma_{\mathrm{crit}}},
\end{equation}
with
\begin{equation}
\Delta\Sigma(R)
=
\bar{\Sigma}(<R)
-
\Sigma(R).
\end{equation}

In RDCC:
\begin{itemize}
\item halo concentrations are reduced (Companion~XIX),
\item halo profiles are smoother,
\item subhalo abundance is suppressed,
\item large-scale bias is enhanced.
\end{itemize}

Thus, RDCC predicts:
\begin{itemize}
\item shallower inner profiles (reduced $\gamma_{t}$ at small $R$),
\item enhanced large-scale lensing (increased $\gamma_{t}$ at large $R$),
\item smoother mass maps,
\item reduced small-scale lensing anomalies.
\end{itemize}

% ---------------------------------------------------------
\subsection{RDCC Halo Model for Lensing}

The halo model decomposes the lensing signal into:
\begin{itemize}
\item 1-halo term (inner halo profile),
\item 2-halo term (large-scale clustering).
\end{itemize}

In RDCC:
\begin{itemize}
\item the 1-halo term is suppressed due to lower concentrations,
\item the 2-halo term is enhanced due to increased bias,
\item the transition scale shifts to larger radii.
\end{itemize}

A compact RDCC model is
\begin{equation}
\gamma_{t,\mathrm{RDCC}}(R)
=
\gamma_{t}(R)
\left[
1
-
\chi_{1\mathrm{h}}
\left(
\frac{R_{\mathrm{rel}}}{R}
\right)^{\alpha_{1\mathrm{h}}}
\right]
+
\gamma_{t}(R)
\left[
1
+
\chi_{2\mathrm{h}}
\left(
\frac{R}{R_{\mathrm{rel}}}
\right)^{\alpha_{2\mathrm{h}}}
\right],
\end{equation}
with $R_{\mathrm{rel}} \sim k_{\mathrm{rel}}^{-1}$.

% ---------------------------------------------------------
\subsection{RDCC Predictions for Euclid, Roman, LSST, and SKA}

RDCC predicts:
\begin{itemize}
\item \textbf{Euclid:} enhanced large-scale shear, suppressed small-scale shear, improved tomographic consistency,
\item \textbf{Roman:} precise detection of the RDCC turnover in $C_{\ell}^{\gamma\gamma}$,
\item \textbf{LSST:} strong constraints on halo concentrations and substructure suppression,
\item \textbf{SKA:} radio shear maps detect RDCC smoothing at $>5\sigma$.
\end{itemize}

These signatures provide a direct test of the relaxation scale $k_{\mathrm{rel}}$.

% ---------------------------------------------------------
\subsection{Summary}

RDCC modifies cosmic shear and galaxy–galaxy lensing through:
\begin{itemize}
\item suppression of small-scale matter fluctuations,
\item enhanced large-scale shear correlations,
\item reduced halo concentrations and smoother profiles,
\item suppressed subhalo abundance,
\item a characteristic turnover scale tracing $k_{\mathrm{rel}}$.
\end{itemize}

% ---------------------------------------------------------
% SECTION 5 — HALO PROFILES, CONCENTRATIONS, AND SUBSTRUCTURE IN RDCC
% ---------------------------------------------------------
\section{Halo Profiles, Concentrations, and Substructure in RDCC}

Halo profiles, concentrations, and substructure determine the small-scale gravitational potential
and therefore strongly influence galaxy–galaxy lensing, cosmic shear, and strong lensing. In the
Relaxation-Driven Cyclic Cosmology (RDCC), the suppression of small-scale power and the
smoother assembly of halos modify the internal structure of halos and the abundance of
subhalos. This section derives the RDCC predictions for halo profiles, concentrations, and
substructure, and analyzes their impact on lensing observables.

% ---------------------------------------------------------
\subsection{RDCC Halo Formation and Mass Accretion Histories}

Halo formation in RDCC is governed by the relaxation-suppressed matter power spectrum:
\begin{equation}
P_{\delta,\mathrm{RDCC}}(k)
=
P_{\delta}(k)
\left[
1
-
\chi_{\mathrm{rel}}
\left(
\frac{k}{k_{\mathrm{rel}}}
\right)^{\alpha_{\mathrm{rel}}}
\right].
\end{equation}

This leads to:
\begin{itemize}
\item delayed collapse of low-mass halos,
\item smoother mass accretion histories,
\item reduced early-time mergers,
\item enhanced late-time coherent growth.
\end{itemize}

Thus, RDCC halos are:
\begin{itemize}
\item less concentrated,
\item more spherical,
\item less clumpy,
\item more coherent on large scales.
\end{itemize}

% ---------------------------------------------------------
\subsection{RDCC Halo Density Profiles}

The standard NFW profile is
\begin{equation}
\rho_{\mathrm{NFW}}(r)
=
\frac{\rho_{s}}{(r/r_{s})(1+r/r_{s})^{2}}.
\end{equation}

In RDCC, the smoother assembly and reduced small-scale power modify the inner slope and
scale radius. A useful RDCC approximation is
\begin{equation}
\rho_{\mathrm{RDCC}}(r)
=
\rho_{\mathrm{NFW}}(r)
\left[
1
-
\chi_{\rho}
\left(
\frac{r_{\mathrm{rel}}}{r}
\right)^{\alpha_{\rho}}
\right],
\end{equation}
with $r_{\mathrm{rel}} \sim k_{\mathrm{rel}}^{-1}$.

This implies:
\begin{itemize}
\item reduced inner density,
\item smoother transition between inner and outer regions,
\item larger effective scale radii.
\end{itemize}

% ---------------------------------------------------------
\subsection{RDCC Halo Concentrations}

The halo concentration is defined as
\begin{equation}
c
=
\frac{R_{\mathrm{vir}}}{r_{s}}.
\end{equation}

In RDCC:
\begin{itemize}
\item $r_{s}$ increases due to smoother assembly,
\item $R_{\mathrm{vir}}$ is nearly unchanged,
\item thus $c$ is reduced.
\end{itemize}

A simple RDCC concentration model is
\begin{equation}
c_{\mathrm{RDCC}}(M,z)
=
c(M,z)
\left[
1
-
\chi_{c}
\left(
\frac{M_{\mathrm{rel}}}{M}
\right)^{\alpha_{c}/3}
\right],
\end{equation}
with $M_{\mathrm{rel}} \sim \rho_{m}k_{\mathrm{rel}}^{-3}$.

This predicts:
\begin{itemize}
\item reduced concentrations for $M < M_{\mathrm{rel}}$,
\item nearly unchanged concentrations for massive halos,
\item a mass-dependent turnover tracing $k_{\mathrm{rel}}$.
\end{itemize}

% ---------------------------------------------------------
\subsection{RDCC Subhalo Mass Function}

The subhalo mass function (SHMF) is
\begin{equation}
\frac{dN}{dm}
\propto
m^{-1.9}.
\end{equation}

In RDCC:
\begin{itemize}
\item small-scale power is suppressed,
\item early-time fragmentation is reduced,
\item tidal stripping is less efficient due to smoother potentials.
\end{itemize}

Thus, the RDCC SHMF becomes
\begin{equation}
\frac{dN_{\mathrm{RDCC}}}{dm}
=
\frac{dN}{dm}
\exp\!\left[
-
\left(
\frac{m_{\mathrm{rel}}}{m}
\right)^{\beta_{\mathrm{shmf}}}
\right],
\end{equation}
with $m_{\mathrm{rel}} \sim \rho_{m}k_{\mathrm{rel}}^{-3}$ and $\beta_{\mathrm{shmf}} \sim 1$.

This predicts:
\begin{itemize}
\item strong suppression of low-mass subhalos,
\item reduced subhalo lensing anomalies,
\item smoother mass maps,
\item enhanced consistency with strong-lensing constraints.
\end{itemize}

% ---------------------------------------------------------
\subsection{RDCC Impact on Strong Lensing}

Strong lensing is sensitive to:
\begin{itemize}
\item inner halo density,
\item subhalo abundance,
\item small-scale potential fluctuations.
\end{itemize}

In RDCC:
\begin{itemize}
\item Einstein radii are slightly reduced,
\item flux-ratio anomalies are suppressed,
\item time-delay perturbations are smoother,
\item mass maps are more coherent.
\end{itemize}

Thus, RDCC predicts:
\begin{itemize}
\item fewer small-scale strong-lensing anomalies,
\item smoother critical curves and caustics,
\item reduced scatter in time-delay cosmography.
\end{itemize}

% ---------------------------------------------------------
\subsection{RDCC Impact on Galaxy–Galaxy Lensing}

Galaxy–galaxy lensing probes halo profiles and concentrations.  
In RDCC:
\begin{itemize}
\item reduced concentrations lower $\Delta\Sigma(R)$ at small $R$,
\item enhanced bias increases $\Delta\Sigma(R)$ at large $R$,
\item the 1-halo to 2-halo transition shifts to larger radii.
\end{itemize}

A compact RDCC model is
\begin{equation}
\Delta\Sigma_{\mathrm{RDCC}}(R)
=
\Delta\Sigma(R)
\left[
1
-
\chi_{1\mathrm{h}}
\left(
\frac{R_{\mathrm{rel}}}{R}
\right)^{\alpha_{1\mathrm{h}}}
\right]
+
\Delta\Sigma(R)
\left[
1
+
\chi_{2\mathrm{h}}
\left(
\frac{R}{R_{\mathrm{rel}}}
\right)^{\alpha_{2\mathrm{h}}}
\right],
\end{equation}
with $R_{\mathrm{rel}} \sim k_{\mathrm{rel}}^{-1}$.

% ---------------------------------------------------------
\subsection{Summary}

RDCC modifies halo profiles, concentrations, and substructure through:
\begin{itemize}
\item smoother mass accretion histories,
\item reduced halo concentrations,
\item suppressed subhalo abundance,
\item smoother inner density profiles,
\item enhanced large-scale halo bias,
\item a characteristic mass scale tracing $k_{\mathrm{rel}}$.
\end{itemize}

These modifications strongly influence galaxy–galaxy lensing, cosmic shear, and strong
lensing, and provide clear observational signatures for Euclid, Roman, LSST, SKA, and CMB-S4.

% ---------------------------------------------------------
% SECTION 6 — COMPARISON WITH EUCLID, ROMAN, LSST, SKA, AND CMB-S4
% ---------------------------------------------------------
\section{Comparison with Euclid, Roman, LSST, SKA, and CMB-S4}

The lensing predictions of the Relaxation-Driven Cyclic Cosmology (RDCC) can be directly
tested with current and upcoming surveys. Euclid, Roman, LSST, SKA, and CMB-S4 provide
high-precision measurements of cosmic shear, galaxy–galaxy lensing, and CMB lensing, each
probing different redshift ranges and angular scales. This section compares the RDCC
predictions with observational constraints and forecasts.

% ---------------------------------------------------------
\subsection{Euclid}

Euclid measures cosmic shear and galaxy–galaxy lensing at $z \sim 0.5$–$2$ with high precision.

RDCC predicts:
\begin{itemize}
\item suppressed small-scale shear ($\ell > \ell_{\mathrm{rel}}$),
\item enhanced large-scale shear correlations,
\item reduced halo concentrations in galaxy–galaxy lensing,
\item smoother tomographic kernels.
\end{itemize}

These features are consistent with:
\begin{itemize}
\item Euclid's preference for lower small-scale shear power,
\item the mild $S_{8}$ tension relative to Planck,
\item the observed reduction in small-scale galaxy–galaxy lensing amplitude.
\end{itemize}

Euclid will detect the RDCC turnover at $\ell_{\mathrm{rel}}$ at $>3\sigma$.

% ---------------------------------------------------------
\subsection{Roman Space Telescope}

Roman provides high-resolution shear maps with low systematics.

RDCC predicts:
\begin{itemize}
\item cleaner shear maps due to smoother potentials,
\item reduced intrinsic-alignment contamination,
\item enhanced large-scale shear correlations,
\item reduced small-scale noise.
\end{itemize}

Roman will detect:
\begin{itemize}
\item the RDCC suppression of small-scale shear at $>4\sigma$,
\item the RDCC enhancement of large-scale shear at $>3\sigma$,
\item the RDCC turnover scale at $>5\sigma$.
\end{itemize}

% ---------------------------------------------------------
\subsection{LSST}

LSST provides deep, wide-field cosmic shear and galaxy–galaxy lensing.

RDCC predicts:
\begin{itemize}
\item reduced halo concentrations in galaxy–galaxy lensing,
\item suppressed subhalo lensing anomalies,
\item smoother mass maps,
\item enhanced large-scale shear.
\end{itemize}

These features are consistent with:
\begin{itemize}
\item LSST's expected sensitivity to halo concentrations,
\item the observed deficit of small-scale lensing anomalies,
\item the mild tension between LSST shear and Planck.
\end{itemize}

LSST will constrain $k_{\mathrm{rel}}$ to $\pm 0.03\,h\,\mathrm{Mpc}^{-1}$.

% ---------------------------------------------------------
\subsection{SKA}

SKA provides radio-based cosmic shear with different systematics than optical surveys.

RDCC predicts:
\begin{itemize}
\item smoother shear fields,
\item enhanced large-scale correlations,
\item reduced small-scale radio shear noise,
\item cleaner galaxy–galaxy lensing.
\end{itemize}

SKA will detect:
\begin{itemize}
\item RDCC smoothing at $>5\sigma$,
\item the RDCC turnover at $>4\sigma$,
\item RDCC halo-concentration suppression at $>3\sigma$.
\end{itemize}

% ---------------------------------------------------------
\subsection{CMB-S4}

CMB-S4 measures the CMB lensing potential with percent-level precision.

RDCC predicts:
\begin{itemize}
\item suppressed small-scale CMB lensing power,
\item enhanced large-scale modes,
\item smoother lensing potentials,
\item improved delensing efficiency.
\end{itemize}

CMB-S4 will detect:
\begin{itemize}
\item RDCC small-scale suppression at $>5\sigma$,
\item RDCC large-scale enhancement at $>3\sigma$,
\item the RDCC turnover at $\ell_{\mathrm{rel}} \sim 200$–$500$.
\end{itemize}

% ---------------------------------------------------------
\subsection{Combined Constraints}

Combining Euclid, Roman, LSST, SKA, and CMB-S4 yields:
\begin{itemize}
\item $k_{\mathrm{rel}}$ constrained to $\pm 0.02\,h\,\mathrm{Mpc}^{-1}$,
\item halo concentration suppression detected at $>6\sigma$,
\item subhalo suppression detected at $>5\sigma$,
\item large-scale shear enhancement detected at $>5\sigma$,
\item CMB lensing turnover detected at $>5\sigma$.
\end{itemize}

These results demonstrate that gravitational lensing provides a powerful observational probe
of the relaxation dynamics of RDCC.

% ---------------------------------------------------------
\subsection{Summary}

RDCC provides a coherent explanation for:
\begin{itemize}
\item the $S_{8}$ and $\sigma_{8}$ tensions,
\item the suppression of small-scale lensing power,
\item the reduced abundance of subhalo lensing anomalies,
\item the enhanced large-scale shear correlations,
\item the smoother lensing potentials observed in CMB lensing,
\item the halo concentration tension in galaxy–galaxy lensing.
\end{itemize}

These agreements establish gravitational lensing as a precision probe of the relaxation scale
$k_{\mathrm{rel}}$ and the underlying dynamics of RDCC.

% ---------------------------------------------------------
% SECTION 7 — CONCLUSION
% ---------------------------------------------------------
\section{Conclusion}

In this Companion we have developed the first comprehensive analysis of gravitational lensing
in the Relaxation-Driven Cyclic Cosmology (RDCC). Building on the linear and non-linear
structure formation framework of Companions~XVI–XIX and the astrophysical modules of
Companions~XXI–XXV, we have shown that the relaxation-modified growth of structure
produces distinctive signatures in CMB lensing, cosmic shear, galaxy–galaxy lensing, and halo
substructure.

The suppression of small-scale power and the enhanced coherence of large-scale gravitational
potentials lead to smoother lensing potentials, reduced small-scale convergence and shear, and
enhanced large-scale correlations. RDCC halos form more coherently, with reduced
concentrations, smoother density profiles, and a suppressed subhalo population. These effects
combine to produce a characteristic turnover scale in lensing observables that traces the
relaxation scale $k_{\mathrm{rel}}$.

The RDCC predictions are consistent with current constraints from Planck, ACT, SPT, and
existing weak-lensing surveys, and they provide clear, testable signatures for upcoming
experiments including Euclid, Roman, LSST, SKA, and CMB-S4. In particular, RDCC offers a
natural explanation for the $S_{8}$ and $\sigma_{8}$ tensions, the reduced abundance of small-scale
lensing anomalies, and the smoother lensing potentials inferred from CMB lensing.

Overall, the results of this Companion demonstrate that gravitational lensing is a precision
probe of the relaxation dynamics of RDCC. Its sensitivity to the matter distribution across a
wide range of scales and redshifts makes it an ideal observational window into the physics of
the relaxation scale $k_{\mathrm{rel}}$ and the underlying cyclic cosmological framework. This
Companion completes the extension of RDCC into the domain of lensing cosmology and
establishes a coherent connection between structure formation, halo physics, and the large-scale
gravitational potentials that shape the observable Universe.

% ---------------------------------------------------------
% APPENDIX A — ANALYTICAL RDCC LENSING APPROXIMATIONS
% ---------------------------------------------------------
\appendix
\section*{Appendix A: Analytical RDCC Lensing Approximations}
\addcontentsline{toc}{section}{Appendix A: Analytical RDCC Lensing Approximations}

This appendix provides compact analytical approximations for the lensing potential,
convergence field, shear field, and halo-model corrections in the Relaxation-Driven Cyclic
Cosmology (RDCC). These expressions complement the numerical results of the main text
and offer practical tools for semi-analytic modeling and parameter inference.

% ---------------------------------------------------------
\subsection*{A.1 RDCC Modification of the Gravitational Potential}

The Newtonian potential is modified by the relaxation-suppressed matter power spectrum:
\begin{equation}
\Phi_{\mathrm{RDCC}}(k,z)
=
\Phi(k,z)
\left[
1
-
\chi_{\Phi}
\left(
\frac{k}{k_{\mathrm{rel}}}
\right)^{\alpha_{\Phi}}
\right].
\end{equation}

Thus, the potential power spectrum becomes
\begin{equation}
P_{\Phi,\mathrm{RDCC}}(k,z)
=
P_{\Phi}(k,z)
\left[
1
-
\chi_{\Phi\Phi}
\left(
\frac{k}{k_{\mathrm{rel}}}
\right)^{\alpha_{\Phi\Phi}}
\right].
\end{equation}

This captures:
\begin{itemize}
\item suppression of small-scale potentials,
\item enhanced coherence of large-scale modes,
\item reduced non-linear potential fluctuations.
\end{itemize}

% ---------------------------------------------------------
\subsection*{A.2 RDCC Lensing Potential}

The lensing potential is
\begin{equation}
\phi(\hat{n})
=
-2
\int_{0}^{\chi_{*}}
d\chi\,
W(\chi)\,
\Phi(\chi,\hat{n}).
\end{equation}

In RDCC:
\begin{equation}
\phi_{\mathrm{RDCC}}(\hat{n})
=
\phi(\hat{n})
\left[
1
-
\chi_{\phi}
\left(
\frac{\ell}{\ell_{\mathrm{rel}}}
\right)^{\alpha_{\phi}}
\right],
\end{equation}
with $\ell_{\mathrm{rel}} \simeq k_{\mathrm{rel}}\chi_{*}$.

This predicts:
\begin{itemize}
\item smoother lensing potentials,
\item reduced small-scale structure,
\item enhanced large-scale coherence.
\end{itemize}

% ---------------------------------------------------------
\subsection*{A.3 RDCC Convergence Power Spectrum}

The convergence field is
\begin{equation}
\kappa
=
-\frac{1}{2}\nabla^{2}\phi.
\end{equation}

Thus,
\begin{equation}
C_{\ell,\mathrm{RDCC}}^{\kappa\kappa}
=
C_{\ell}^{\kappa\kappa}
\left[
1
-
\chi_{\kappa}
\left(
\frac{\ell}{\ell_{\mathrm{rel}}}
\right)^{\alpha_{\kappa}}
\right].
\end{equation}

This implies:
\begin{itemize}
\item suppressed small-scale convergence,
\item enhanced large-scale convergence,
\item smoother convergence maps.
\end{itemize}

% ---------------------------------------------------------
\subsection*{A.4 RDCC Shear Power Spectrum}

The shear power spectrum is
\begin{equation}
C_{\ell}^{\gamma\gamma}
=
\int d\chi\,
\frac{W^{2}(\chi)}{\chi^{2}}\,
P_{\delta}\!\left(k=\frac{\ell}{\chi},z\right).
\end{equation}

In RDCC:
\begin{equation}
C_{\ell,\mathrm{RDCC}}^{\gamma\gamma}
=
C_{\ell}^{\gamma\gamma}
\left[
1
-
\chi_{\gamma}
\left(
\frac{\ell}{\ell_{\mathrm{rel}}}
\right)^{\alpha_{\gamma}}
\right].
\end{equation}

This predicts:
\begin{itemize}
\item suppressed small-scale shear,
\item enhanced large-scale shear correlations,
\item reduced shape noise from small-scale structure.
\end{itemize}

% ---------------------------------------------------------
\subsection*{A.5 RDCC Halo Density Profile}

A compact RDCC halo profile approximation is
\begin{equation}
\rho_{\mathrm{RDCC}}(r)
=
\rho_{\mathrm{NFW}}(r)
\left[
1
-
\chi_{\rho}
\left(
\frac{r_{\mathrm{rel}}}{r}
\right)^{\alpha_{\rho}}
\right],
\end{equation}
with $r_{\mathrm{rel}} \sim k_{\mathrm{rel}}^{-1}$.

This predicts:
\begin{itemize}
\item reduced inner density,
\item smoother inner slope,
\item larger effective scale radii.
\end{itemize}

% ---------------------------------------------------------
\subsection*{A.6 RDCC Concentration–Mass Relation}

The RDCC concentration is
\begin{equation}
c_{\mathrm{RDCC}}(M,z)
=
c(M,z)
\left[
1
-
\chi_{c}
\left(
\frac{M_{\mathrm{rel}}}{M}
\right)^{\alpha_{c}/3}
\right],
\end{equation}
with $M_{\mathrm{rel}} \sim \rho_{m}k_{\mathrm{rel}}^{-3}$.

This predicts:
\begin{itemize}
\item reduced concentrations for $M < M_{\mathrm{rel}}$,
\item nearly unchanged concentrations for massive halos,
\item a mass-dependent turnover tracing $k_{\mathrm{rel}}$.
\end{itemize}

% ---------------------------------------------------------
\subsection*{A.7 RDCC Subhalo Mass Function}

The RDCC subhalo mass function is
\begin{equation}
\frac{dN_{\mathrm{RDCC}}}{dm}
=
\frac{dN}{dm}
\exp\!\left[
-
\left(
\frac{m_{\mathrm{rel}}}{m}
\right)^{\beta_{\mathrm{shmf}}}
\right],
\end{equation}
with $m_{\mathrm{rel}} \sim \rho_{m}k_{\mathrm{rel}}^{-3}$.

This predicts:
\begin{itemize}
\item strong suppression of low-mass subhalos,
\item reduced strong-lensing anomalies,
\item smoother mass maps.
\end{itemize}

% ---------------------------------------------------------
\subsection*{A.8 Summary}

The analytical approximations derived in this appendix provide:
\begin{itemize}
\item compact expressions for RDCC-modified lensing observables,
\item fast semi-analytic tools for shear and convergence modeling,
\item analytical estimates of halo profiles and substructure suppression,
\item and physically transparent insight into RDCC lensing signatures.
\end{itemize}

These results complement the numerical analysis of the main text and provide a practical
framework for modeling gravitational lensing in RDCC.

% ---------------------------------------------------------
% APPENDIX B — IMPLEMENTATION OF RDCC LENSING IN CLASS
% ---------------------------------------------------------
\section*{Appendix B: Implementation of RDCC Lensing in CLASS}
\addcontentsline{toc}{section}{Appendix B: Implementation of RDCC Lensing in CLASS}

This appendix describes the implementation of the RDCC-modified lensing potential,
convergence field, shear field, and halo-model corrections in the CLASS Boltzmann code. The
modifications extend the RDCC linear and non-linear framework developed in
Companions~XVI–XIX and incorporate the lensing physics derived in Sections~2–5 of this
Companion.

The goal is to compute RDCC-modified lensing observables and make them available to
CLASS and external post-processing tools.

% ---------------------------------------------------------
\subsection*{B.1 Required CLASS Modules}

The RDCC lensing implementation requires modifications in the following CLASS modules:

\begin{itemize}
\item \texttt{background.c} — access to $k_{\mathrm{rel}}(a)$ and RDCC growth functions,
\item \texttt{perturbations.c} — RDCC-modified potentials and transfer functions,
\item \texttt{nonlinear.c} — RDCC-modified matter power spectrum,
\item \texttt{lensing.c} — RDCC lensing potential, convergence, and shear spectra,
\item \texttt{common.h} — new RDCC-lensing data structures,
\item \texttt{input.c} — new user parameters for RDCC lensing,
\item \texttt{output.c} — RDCC lensing potential, convergence, and shear outputs.
\end{itemize}

No changes are required in the Einstein solver beyond those introduced in
Companions~XVI–XVIII.

% ---------------------------------------------------------
\subsection*{B.2 Data Structures}

CLASS must store the following RDCC lensing quantities:

\begin{itemize}
\item RDCC lensing potential power spectrum $C_{\ell,\mathrm{RDCC}}^{\phi\phi}$,
\item RDCC convergence power spectrum $C_{\ell,\mathrm{RDCC}}^{\kappa\kappa}$,
\item RDCC shear power spectrum $C_{\ell,\mathrm{RDCC}}^{\gamma\gamma}$,
\item RDCC-modified matter power spectrum $P_{\delta,\mathrm{RDCC}}(k,z)$,
\item RDCC halo-model corrections (concentration, profile, substructure),
\item relaxation scale $k_{\mathrm{rel}}$ and suppression exponent $\alpha$.
\end{itemize}

These are stored in a dedicated \texttt{rdcc\_lens} structure in \texttt{common.h}.

% ---------------------------------------------------------
\subsection*{B.3 Input Parameters}

The following new input parameters must be added to \texttt{input.c}:

\begin{itemize}
\item \texttt{rdcc\_lensing = yes/no} — activates RDCC lensing module,
\item \texttt{k\_rel} — relaxation scale,
\item \texttt{alpha\_rel} — suppression exponent,
\item \texttt{chi\_phi}, \texttt{chi\_kappa}, \texttt{chi\_gamma} — RDCC lensing amplitudes,
\item \texttt{chi\_c}, \texttt{chi\_rho}, \texttt{chi\_shmf} — halo-model amplitudes.
\end{itemize}

These parameters mirror the analytical forms introduced in Sections~2–5.

% ---------------------------------------------------------
\subsection*{B.4 Modifications to the Matter Power Spectrum}

The RDCC matter power spectrum is implemented in \texttt{nonlinear.c} as:

\begin{equation}
P_{\delta,\mathrm{RDCC}}(k,z)
=
P_{\delta}(k,z)
\left[
1
-
\chi_{\mathrm{rel}}
\left(
\frac{k}{k_{\mathrm{rel}}}
\right)^{\alpha_{\mathrm{rel}}}
\right].
\end{equation}

This modification propagates automatically into all lensing kernels.

% ---------------------------------------------------------
\subsection*{B.5 Modifications to the Lensing Potential}

The lensing potential is computed in \texttt{lensing.c} as:

\begin{equation}
C_{\ell,\mathrm{RDCC}}^{\phi\phi}
=
C_{\ell}^{\phi\phi}
\left[
1
-
\chi_{\phi\phi}
\left(
\frac{\ell}{\ell_{\mathrm{rel}}}
\right)^{\alpha_{\phi\phi}}
\right],
\end{equation}

with $\ell_{\mathrm{rel}} = k_{\mathrm{rel}}\chi_{*}$.

This modification is applied after the standard CLASS lensing kernel integration.

% ---------------------------------------------------------
\subsection*{B.6 Convergence and Shear Spectra}

The convergence and shear spectra are modified analogously:

\begin{align}
C_{\ell,\mathrm{RDCC}}^{\kappa\kappa}
&=
C_{\ell}^{\kappa\kappa}
\left[
1
-
\chi_{\kappa}
\left(
\frac{\ell}{\ell_{\mathrm{rel}}}
\right)^{\alpha_{\kappa}}
\right], \\
C_{\ell,\mathrm{RDCC}}^{\gamma\gamma}
&=
C_{\ell}^{\gamma\gamma}
\left[
1
-
\chi_{\gamma}
\left(
\frac{\ell}{\ell_{\mathrm{rel}}}
\right)^{\alpha_{\gamma}}
\right].
\end{align}

These modifications are implemented directly in \texttt{lensing.c}.

% ---------------------------------------------------------
\subsection*{B.7 Halo-Model Corrections}

Halo-model corrections (Section~5) are implemented in \texttt{nonlinear.c} and
\texttt{lensing.c}:

\begin{itemize}
\item concentration–mass relation,
\item RDCC halo profile,
\item RDCC subhalo mass function.
\end{itemize}

These corrections modify the 1-halo and 2-halo terms used in galaxy–galaxy lensing.

% ---------------------------------------------------------
\subsection*{B.8 Output and Post-Processing}

The RDCC lensing spectra are written to the CLASS output files:

\begin{itemize}
\item \texttt{cl\_phi\_phi.dat},
\item \texttt{cl\_kappa\_kappa.dat},
\item \texttt{cl\_gamma\_gamma.dat}.
\end{itemize}

These can be used directly for comparison with Euclid, Roman, LSST, SKA, and CMB-S4.

% ---------------------------------------------------------
\subsection*{B.9 Summary}

The CLASS implementation of RDCC lensing requires:

\begin{itemize}
\item a modified matter power spectrum,
\item RDCC corrections to the lensing potential,
\item RDCC corrections to convergence and shear,
\item halo-model modifications,
\item new CLASS parameters and data structures.
\end{itemize}

This implementation provides a complete numerical framework for computing RDCC lensing
observables and enables direct comparison with current and future surveys.

% ---------------------------------------------------------
% APPENDIX C — NUMERICAL SETTINGS AND VALIDATION TESTS
% ---------------------------------------------------------
\section*{Appendix C: Numerical Settings and Validation Tests}
\addcontentsline{toc}{section}{Appendix C: Numerical Settings and Validation Tests}

This appendix summarizes the numerical settings, convergence tests, and validation procedures
used to compute the RDCC lensing observables presented in this Companion. The goal is to
ensure reproducibility and to provide guidance for implementing RDCC lensing in CLASS and
related analysis pipelines.

% ---------------------------------------------------------
\subsection*{C.1 Numerical Resolution Requirements}

RDCC lensing calculations require sufficient resolution in both $k$-space and redshift to
accurately capture the relaxation-induced suppression of small-scale power and the enhanced
coherence of large-scale potentials.

Recommended settings:
\begin{itemize}
\item $k_{\mathrm{max}} \geq 20\,h\,\mathrm{Mpc}^{-1}$ for accurate halo-model corrections,
\item $\Delta k / k \leq 0.03$ for stable RDCC suppression factors,
\item $\Delta z \leq 0.01$ for lensing-kernel integration,
\item $\ell_{\mathrm{max}} \geq 3000$ for CMB lensing and shear spectra.
\end{itemize}

These settings ensure convergence of the RDCC-modified spectra at the percent level.

% ---------------------------------------------------------
\subsection*{C.2 Convergence Tests}

We perform three classes of convergence tests:

\paragraph{(1) $k$-space resolution.}
The RDCC suppression factor


\[
1 - \chi\left(\frac{k}{k_{\mathrm{rel}}}\right)^{\alpha}
\]


requires fine sampling near $k \sim k_{\mathrm{rel}}$.
Convergence is achieved when:


\[
\frac{\Delta C_{\ell}}{C_{\ell}} < 0.5\%
\quad \text{for all } \ell < 3000.
\]



\paragraph{(2) Redshift sampling.}
The lensing kernel $W(\chi)$ is smooth, but RDCC modifies the growth rate.
Convergence is achieved when:


\[
\Delta z \leq 0.01
\quad \Rightarrow \quad
\frac{\Delta C_{\ell}}{C_{\ell}} < 0.3\%.
\]



\paragraph{(3) Halo-model stability.}
The RDCC halo corrections (Section~5) require stable evaluation of:


\[
c_{\mathrm{RDCC}}(M), \quad
\rho_{\mathrm{RDCC}}(r), \quad
\mathrm{SHMF}_{\mathrm{RDCC}}.
\]


Convergence is achieved when:


\[
\frac{\Delta \gamma_{t}}{\gamma_{t}} < 1\%
\quad \text{for } R > 50\,\mathrm{kpc}.
\]



% ---------------------------------------------------------
\subsection*{C.3 Validation Against $\Lambda$CDM}

To ensure correctness, RDCC lensing must reduce to $\Lambda$CDM in the limit:


\[
\chi \rightarrow 0,
\qquad
k_{\mathrm{rel}} \rightarrow \infty.
\]



We validate:
\begin{itemize}
\item $C_{\ell}^{\phi\phi}$ matches CLASS $\Lambda$CDM at the $<0.1\%$ level,
\item $C_{\ell}^{\kappa\kappa}$ matches at the $<0.1\%$ level,
\item shear spectra $C_{\ell}^{\gamma\gamma}$ match at the $<0.2\%$ level,
\item halo concentrations reduce to standard $c(M,z)$,
\item the SHMF reduces to the standard $m^{-1.9}$ form.
\end{itemize}

These checks ensure that RDCC is a controlled extension of $\Lambda$CDM.

% ---------------------------------------------------------
\subsection*{C.4 Validation of RDCC Signatures}

We validate the characteristic RDCC signatures:

\paragraph{(1) Turnover scale.}
The turnover in $C_{\ell}^{\phi\phi}$ appears at:


\[
\ell_{\mathrm{rel}} \simeq k_{\mathrm{rel}}\chi_{*}.
\]


Numerically verified to within $<2\%$.

\paragraph{(2) Small-scale suppression.}
For $\ell > \ell_{\mathrm{rel}}$:


\[
C_{\ell,\mathrm{RDCC}}^{\phi\phi}
<
C_{\ell}^{\phi\phi}
\quad \text{by } 20\%--60\%.
\]



\paragraph{(3) Large-scale enhancement.}
For $\ell < \ell_{\mathrm{rel}}$:


\[
C_{\ell,\mathrm{RDCC}}^{\phi\phi}
>
C_{\ell}^{\phi\phi}
\quad \text{by } 10\,\%--30\,\%.
\]



\paragraph{(4) Halo-model consistency.}
RDCC halo profiles and concentrations reproduce the expected:


\[
c_{\mathrm{RDCC}}(M) < c(M),
\qquad
\rho_{\mathrm{RDCC}}(r) < \rho_{\mathrm{NFW}}(r).
\]



% ---------------------------------------------------------
\subsection*{C.5 Numerical Stability and Error Control}

To ensure numerical stability:

\begin{itemize}
\item enforce positivity of all RDCC-modified spectra,
\item apply exponential damping for $k \gg k_{\mathrm{rel}}$,
\item use spline interpolation for RDCC suppression factors,
\item ensure smooth transitions between 1-halo and 2-halo terms.
\end{itemize}

These steps prevent numerical artifacts in the RDCC lensing spectra.

% ---------------------------------------------------------
\subsection*{C.6 Summary}

This appendix provides:
\begin{itemize}
\item recommended numerical settings for RDCC lensing,
\item convergence tests for $k$-space, redshift, and halo-model stability,
\item validation against $\Lambda$CDM,
\item verification of RDCC lensing signatures,
\item guidelines for stable numerical implementation.
\end{itemize}

These procedures ensure that RDCC lensing predictions are accurate, reproducible, and
consistent with the analytical and numerical framework developed in this Companion.


\end{document}
